% ==============================================================================
%  6.1  El modelo de regresión lineal múltiple
% ==============================================================================
\section{El modelo de regresión lineal múltiple}
\label{sec:t6-modelo}

Se quiere modelizar la esperanza de una variable continua $Y$ cuando se conocen $p-1$
variables explicativas (predictores o regresores) $X_1,X_2,\ldots,X_{p-1}$. El modelo es
``lineal'' porque la esperanza de $Y$ condicionada a las $X$ depende linealmente de
$X_1,X_2,\ldots,X_{p-1}$.

Por ejemplo, con dos predictores ($p=3$) se propone el modelo
\[
  \E(Y\cond X_1,X_2) = \beta_0+\beta_1X_1+\beta_2X_2 ,
\]
donde:
\begin{itemize}
  \item $\beta_0$, $\beta_1$ y $\beta_2$ son constantes desconocidas, llamadas
    \textbf{parámetros del modelo} o \textbf{coeficientes de la regresión};
  \item se estiman a partir de una muestra aleatoria de $n$ observaciones $(X_{i1},X_{i2},Y_i)$,
    $i=1,\ldots,n$, donde $Y_i$ es la respuesta medida en el $i$-ésimo individuo y $X_{i1}$,
    $X_{i2}$ son los valores de las variables explicativas en ese individuo.
\end{itemize}
Por simplicidad, se escribe $\E(Y)$ en lugar de $\E(Y\cond X_1,X_2)$.

\begin{definicion}[Modelo de regresión lineal múltiple][def:t6-rlm]
  \[
    Y_i = \beta_0+\beta_1X_{i1}+\beta_2X_{i2}+\cdots+\beta_{p-1}X_{i,p-1}+\varepsilon_i,
    \qquad \varepsilon_i\iid\Normal(0,\sigma^2), \quad i=1,\ldots,n .
  \]
  La función $\E(Y\cond X_1,X_2,\ldots,X_{p-1})$ se llama \textbf{función de respuesta}.
\end{definicion}

En el modelo de regresión lineal simple, $\E(Y\cond X_1)$ es una recta; con dos variables
explicativas, la función de respuesta $\E(Y\cond X_1,X_2)$ es un plano. A partir de $p>3$ ya no es
posible representarla gráficamente. En regresión múltiple, la función de respuesta también se
denomina \textbf{superficie de regresión} o \textbf{superficie de respuesta}.

\begin{figure}[H]
  \centering
  \begin{tikzpicture}
    \begin{axis}[view={50}{22}, width=10cm, height=8cm, xmin=0, xmax=10, ymin=0, ymax=10,
        zmin=0, zmax=90, xlabel={$X_1$}, ylabel={$X_2$}, zlabel={$Y$}, ticks=none,
        axis lines=left, clip=false, label style={font=\small}]
      \addplot3[surf, shader=flat, fill=NavyBlue!12, faceted color=NavyBlue!60, opacity=0.85,
        domain=0:10, y domain=0:10, samples=5, samples y=5] {10+2*x+5*y};
      % observación (X_i1, X_i2) = (6, 4)
      \draw[dashed, gray] (axis cs:6,4,0) -- (axis cs:6,4,42);
      \draw[dashed, gray] (axis cs:6,0,0) -- (axis cs:6,4,0) -- (axis cs:0,4,0);
      \fill (axis cs:6,4,0) circle (1.6pt);
      \node[font=\scriptsize, anchor=north west] at (axis cs:6,4,0) {$(X_{i1},X_{i2})$};
      \fill[NavyBlue] (axis cs:6,4,42) circle (1.8pt);
      \node[font=\scriptsize, NavyBlue, anchor=west] at (axis cs:6.3,4,40) {$\E(Y_i)$};
      \draw[colRes, thick] (axis cs:6,4,42) -- (axis cs:6,4,56);
      \fill[colRes] (axis cs:6,4,56) circle (1.8pt);
      \node[font=\scriptsize, colRes, anchor=west] at (axis cs:6.3,4,57) {$Y_i$};
      \node[font=\scriptsize, colRes, anchor=east] at (axis cs:5.8,4,49) {$\varepsilon_i$};
      \node[font=\scriptsize, anchor=east] at (axis cs:0,0,10) {$\beta_0=10$};
    \end{axis}
  \end{tikzpicture}
  \caption{Plano de respuesta $\E(Y)=10+2X_1+5X_2$. Para la observación $i$, la respuesta
    $Y_i$ se desvía de su valor esperado $\E(Y_i)$, situado sobre el plano, en el error
    $\varepsilon_i$.}
\end{figure}

\subsection{Hipótesis del modelo}

Los supuestos sobre la componente aleatoria son los mismos que en la regresión simple
(\cref{sec:t2-hipotesis}): $\varepsilon_1,\ldots,\varepsilon_n$ son variables aleatorias
i.i.d.\ con $\varepsilon_i\sim\Normal(0,\sigma^2)$. Es decir:
\begin{itemize}
  \item los $\varepsilon_i$ tienen media cero, $\E(\varepsilon_i)=0$ (\textbf{linealidad});
  \item los $\varepsilon_i$ tienen todos la misma varianza $\sigma^2$, normalmente desconocida,
    que es otro parámetro del modelo a estimar (\textbf{homogeneidad de la varianza});
  \item los $\varepsilon_i$ tienen distribución normal (\textbf{normalidad});
  \item los $\varepsilon_i$ son independientes entre sí e independientes de los predictores
    $x_{i1},x_{i2},\ldots,x_{i,p-1}$ (\textbf{independencia}).
\end{itemize}
Estas hipótesis equivalen a que, conocidos $X_{i1},X_{i2},\ldots,X_{i,p-1}$,
\[
  Y_i\sim\Normal\left(\beta_0+\sum_{j=1}^{p-1}\beta_jX_{ij},\ \sigma^2\right)\ \text{independientes}.
\]

\subsection{Interpretación de los coeficientes}

En el caso con dos predictores, $\E(Y)=\beta_0+\beta_1X_1+\beta_2X_2$:
\begin{itemize}
  \item $\beta_0$ es el \textbf{\emph{intercept}} u ordenada en el origen del plano: la respuesta media
    en el punto $(X_1,X_2)=(0,0)$;
  \item $\beta_1$ indica el cambio en la respuesta media cuando $X_1$ aumenta una unidad,
    manteniendo $X_2$ constante;
  \item análogamente, $\beta_2$ indica el cambio en la respuesta media cuando $X_2$ aumenta una
    unidad, manteniendo $X_1$ constante.
\end{itemize}

\begin{ejemplo}[Interpretación de los coeficientes][ej:t6-interpretacion]
  Sea $\E(Y)=10+2X_1+5X_2$. Si se fija $X_2=3$, se obtiene la recta $\E(Y)=25+2X_1$. Lo mismo
  ocurre para cualquier otro valor de $X_2$: solo cambia la constante de la recta, $10+5X_2$.
  Por tanto, $\beta_1=2$ indica que la respuesta media aumenta 2 unidades por cada incremento
  unitario de $X_1$ cuando $X_2$ se mantiene constante, sea cual sea el valor de $X_2$.
\end{ejemplo}

Por ello, $\beta_1$ y $\beta_2$ se llaman a veces \textbf{coeficientes de regresión parcial}:
reflejan el efecto parcial de una variable explicativa cuando el otro predictor está incluido en el
modelo y se mantiene constante.

\subsection{Consideraciones sobre el modelo}

\begin{itemize}
  \item El término \textbf{``lineal''} se refiere a los \textbf{parámetros}, no a las variables
    explicativas (\cref{def:modelo-lineal}):
    \begin{itemize}
      \item pueden utilizarse modelos de regresión lineales para manejar cualquier función de
        una variable explicativa, por ejemplo $X^2$ o $\log(X)$: $\beta_0+\beta_1X+\beta_2X^2$;
      \item no se utilizan modelos de regresión lineales para funciones no lineales de los
        parámetros, como $\beta_0+\beta_1X_1+\beta_1^2X_2$.
    \end{itemize}
  \item Los \textbf{regresores} pueden ser:
    \begin{itemize}
      \item variables continuas o cualitativas;
      \item variables transformadas;
      \item potencias de una misma variable (\textbf{regresión polinomial}):
        $Y_i=\beta_0+\beta_1X_i+\beta_2X_i^2+\cdots+\beta_{p-1}X_i^{p-1}+\varepsilon_i$;
      \item \textbf{efectos de interacción}: el producto de predictores se incluye como una
        variable adicional, $Y_i=\beta_0+\beta_1X_{i1}+\beta_2X_{i2}+\beta_3X_{i1}X_{i2}+\varepsilon_i$.
    \end{itemize}
  \item \textbf{Multicolinealidad}: si las variables explicativas están muy correlacionadas entre
    sí, es posible que no se obtengan buenos estimadores de los parámetros, ya que tratan de
    explicar la misma parte de la variabilidad de la respuesta (\cref{sec:t6-multicolinealidad}).
\end{itemize}

\subsection{El modelo en notación matricial}
\label{sec:t6-matricial}

El modelo múltiple se escribe como en el Tema 5 (\cref{def:t5-modelo}),
$\vect{Y}=\vect{X}\vect{\beta}+\vect{\varepsilon}$ con
$\vect{\varepsilon}\sim\Normal(\vect{0},\sigma^2\vect{I}_n)$, añadiendo una columna a la matriz de
diseño por cada predictor:
\[
  \vect{X}_{n\times p} =
  \begin{pmatrix}
    1 & X_{11} & X_{12} & \cdots & X_{1,p-1}\\
    1 & X_{21} & X_{22} & \cdots & X_{2,p-1}\\
    \vdots & \vdots & \vdots & & \vdots\\
    1 & X_{n1} & X_{n2} & \cdots & X_{n,p-1}
  \end{pmatrix},
  \qquad
  \vect{\beta}_{p\times1} = \begin{pmatrix} \beta_0\\ \beta_1\\ \vdots\\ \beta_{p-1} \end{pmatrix}.
\]
Todos los resultados del Tema 5 se mantienen con esta $\vect{X}$: el estimador
$\bhat=\XtXi\vect{X}\T\vect{Y}$ (si $\XtX$ es invertible), la matriz hat, las sumas de cuadrados y
$\bhat\sim\Normal\big(\vect{\beta},\sigma^2\XtXi\big)$. Lo único que cambia son los grados de
libertad: como se estiman $p$ coeficientes, $SSE$ tiene $n-p$ grados de libertad, el estimador
insesgado de $\sigma^2$ es
\[
  MSE = \frac{SSE}{n-p},
\]
y los estadísticos $t$ para los coeficientes, la respuesta media y la predicción
(\cref{sec:t5-respuesta-media}) siguen una $\tdist{n-p}$. $SSR$ tiene $p-1$ grados de libertad,
uno por predictor.

\begin{center}
  \renewcommand{\arraystretch}{1.3}
  \begin{tabular}{lccccc}
    \toprule
    Fuente & $SS$ & g.l. & $MS$ & $F$ & $p$-valor \\
    \midrule
    Modelo (regresión) & $SSR$ & $p-1$ & $MSR=SSR/(p-1)$ & $F^*=MSR/MSE$ & $p$ \\
    Error & $SSE$ & $n-p$ & $MSE=SSE/(n-p)$ & & \\
    Total & $SST$ & $n-1$ & & & \\
    \bottomrule
  \end{tabular}
\end{center}
El estadístico $F=MSR/MSE$ contrasta globalmente $H_0\colon\beta_1=\cdots=\beta_{p-1}=0$ frente a
que algún $\beta_k$ sea distinto de 0, con distribución $\Fdist{p-1}{n-p}$ bajo $H_0$; es el
contraste de la regresión (\cref{sec:t3-anova}) con $p-1$ predictores. Cada coeficiente se
contrasta individualmente con el estadístico $t$ de $H_0\colon\beta_k=0$.

\begin{nota}[Sumas de cuadrados de tipo I y de tipo III]
  Las salidas de software descomponen $SSR$ en la aportación de cada predictor, como reducción de
  $SSE$ al incluirlo en el modelo (\cref{sec:t3-ftest-general}):
  \begin{itemize}
    \item las \textbf{sumas de cuadrados de tipo I} (secuenciales) miden la reducción de $SSE$ al
      añadir cada variable al modelo que contiene las anteriores, en el orden en que se
      introducen; suman $SSR$;
    \item las \textbf{sumas de cuadrados de tipo III} (parciales) miden la reducción de $SSE$ al
      añadir cada variable al modelo que contiene todas las demás.
  \end{itemize}
\end{nota}
